% Part: second-order-logic
% Chapter: syntax-and-semantics
% Section: size-of-domain

\documentclass[../../../include/open-logic-section]{subfiles}

\begin{document}

\olfileid[pt]{sol}{syn}{siz}

\olsection{Infinitude e \usetoken{S}{enumerable}
  \usetoken{P}{domain}}

Um conjunto~$M$ é infinito (no sentido de Dedekind) sse existe
!!a{injective} função~$f\colon M \to M$ que não é !!{surjective}, isto é,
com $\ran{f} \neq M$. Na lógica de primeira ordem, podemos considerar uma
!!{function} unária~$f$ e dizer que a função~$\Assign{f}{M}$ atribuída a ela
em !!a{structure}~$\Struct{M}$ é !!{injective} e
$\ran{f} \neq \Domain{M}$:
\[
\lforall[x][\lforall[y][(\eq[f(x)][f(y)] \to \eq[x][y])]] \land
\lexists[y][\lforall[x][\eq/[y][f(x)]]].
\]
Se $\Struct{M}$ satisfaz essa !!{sentence}, então
$\Assign{f}{M}: \Domain{M} \to \Domain{M}$ é !!{injective}, e portanto
$\Domain{M}$ deve ser infinito. Se $\Domain{M}$ é infinito e, assim, existe
tal função, podemos tomar $\Assign{f}{M}$ como essa função, e $\Struct{M}$
satisfará a !!{sentence}. Contudo, isso exige que nossa linguagem contenha o
símbolo não lógico~$f$ usado para esse fim. Na lógica de segunda ordem,
podemos simplesmente dizer que tal função \emph{existe}. Isso já não exige
~$f$, e obtemos, na lógica pura de segunda ordem, a !!{sentence}
\[
\fn{Inf} \ident \lexists[u][(\lforall[x][\lforall[y][(\eq[u(x)][u(y)]
      \lif \eq[x][y])]] \land \lexists[y][\lforall[x][\eq/[y][u(x)]]])].
\]
$\Sat{M}{\fn{Inf}}$ sse $\Domain{M}$ é infinito. Podemos então definir
$\fn{Fin} \ident \lnot \fn{Inf}$; $\Sat{M}{\fn{Fin}}$ sse $\Domain{M}$ é
finito. Nenhuma !!{sentence} isolada da lógica pura de primeira ordem pode
expressar que o !!{domain} é infinito, embora um conjunto infinito delas
possa. Não existe conjunto de !!{sentence}s da lógica pura de primeira
ordem que seja satisfeito em !!a{structure} sse seu domínio é finito.

\begin{prop}
$\Sat{M}{\fn{Inf}}$ sse $\Domain{M}$ é infinito.
\end{prop}

\begin{proof}
$\Sat{M}{\fn{Inf}}$ sse
  $\Sat{M}{\lforall[x][\lforall[y][(\eq[u(x)][u(y)] \lif \eq[x][y])]]
    \land \lexists[y][\lforall[x][\eq/[y][u(x)]]]}[s]$ para algum~$s$.
  Se isso vale, $s(u)$ é !!a{injective} função, e algum
  $y \in \Domain{M}$ não pertence à imagem de~$s(u)$. Reciprocamente, se
  existe !!a{injective} $f\colon \Domain{M} \to \Domain{M}$ com
  $\ran{f} \neq \Domain{M}$, então $s(u) = f$ é uma atribuição de variáveis
  desse tipo.
\end{proof}

Um conjunto $M$ é !!{enumerable} se existe uma enumeração
\[
m_0, m_1, m_2, \dots
\]
de seus !!{element}s (sem repetições, mas possivelmente finita). Tal
enumeração existe sse há !!a{element} $z \in M$ e uma função
$f\colon M \to M$ tais que $z$, $f(z)$, $f(f(z))$, \dots, são todos os
!!{element}s de~$M$. Pois, se existe a enumeração, $z = m_0$ e
$f(m_k) = m_{k+1}$ (ou $f(m_k) = m_k$ se $m_k$ é o último !!{element} da
enumeração) são o !!{element} e a função requeridos. Por outro lado, se
existem tais $z$ e $f$, então $z$, $f(z)$, $f(f(z))$, \dots, é uma
enumeração de~$M$, e $M$ é !!{enumerable}. Podemos expressar a existência de
$z$ e~$f$ na lógica de segunda ordem para produzir uma !!{sentence}
verdadeira em !!a{structure} sse a !!{structure} é !!{enumerable}:
\[
\fn{Count} \ident
\lexists[z][\lexists[u][\lforall[X][((X(z) \land
      \lforall[x][(X(x) \lif X(u(x)))]) \lif \lforall[x][X(x)])]]]
\]

\begin{prop}
$\Sat{M}{\fn{Count}}$ sse $\Domain{M}$ é !!{enumerable}.
\end{prop}

\begin{proof}
Suponha que $\Domain{M}$ seja !!{enumerable}, e sejam $m_0$, $m_1$, \dots,
uma enumeração. Removendo repetições, podemos garantir que nenhum $m_k$
apareça duas vezes. Defina $f(m_k) = m_{k+1}$ e tome $s(z) = m_0$ e
$s(u) = f$. Mostramos que
\[
\Sat{M}{\lforall[X][((X(z) \land \lforall[x][(X(x) \lif X(u(x)))])
    \lif \lforall[x][X(x)])]}[s]
\]
Suponha que $M \subseteq \Domain{M}$ seja arbitrário. Suponha ainda que
$\Sat{M}{(X(z) \land \lforall[x][(X(x) \lif
X(u(x)))])}[\Subst{s}{M}{X}]$. Então $\Subst{s}{M}{X}(z) \in M$, e sempre
que $x \in M$, também $(\Subst{s}{M}{X}(u))(x) \in M$. Em outras palavras,
como $\varAssign{\Subst{s}{M}{X}}{s}{X}$, $m_0 \in M$ e, se $x \in M$,
então $f(x) \in M$; portanto, $m_0 \in M$, $m_1 = f(m_0) \in M$,
$m_2 = f(f(m_0)) \in M$ etc. Assim, $M = \Domain{M}$ e, portanto,
$\Sat{M}{\lforall[x][X(x)]}[\Subst{s}{M}{X}]$. Como
$M \subseteq \Domain{M}$ era arbitrário, concluímos:
$\Sat{M}{\fn{Count}}$.

Suponha agora que $\Sat{M}{\fn{Count}}$, isto é,
\[
\Sat{M}{\lforall[X][((X(z) \land \lforall[x][(X(x) \lif X(u(x)))])
    \lif \lforall[x][X(x)])]}[s]
\]
para algum~$s$. Sejam $m = s(z)$ e $f = s(u)$, e considere
$M = \{m, f(m), f(f(m)), \dots\}$. O conjunto $M$ assim definido é
claramente !!{enumerable}. Então
\[
\Sat{M}{(X(z) \land \lforall[x][(X(x) \lif X(u(x)))])
    \lif \lforall[x][X(x)]}[\Subst{s}{M}{X}]
\]
por hipótese. Além disso, $\Sat{M}{X(z)}[\Subst{s}{M}{X}]$, pois $M \ni m =
\Subst{s}{M}{X}(z)$, e também
$\Sat{M}{\lforall[x][(X(x) \lif X(u(x)))]}[\Subst{s}{M}{X}]$, pois, sempre
que $x \in M$, também $f(x) \in M$. Portanto, como tanto o antecedente
quanto a condicional são satisfeitos, o consequente também deve ser:
$\Sat{M}{\lforall[x][X(x)]}[\Subst{s}{M}{X}]$. Mas isso significa que
$M = \Domain{M}$ e, assim, $\Domain{M}$ é !!{enumerable}, pois $M$ o é por
definição.
\end{proof}

\begin{prob}
A !!{sentence} $\fn{Inf} \land \fn{Count}$ é verdadeira precisamente nos
!!{denumerable} domínios. Ajuste a definição de $\fn{Count}$ para que ela se
torne uma !!{sentence} diferente que expresse diretamente que o domínio é
!!{denumerable}, e prove que o faz.
\end{prob}

\end{document}
